\section{Small-set returns beyond a fixed-power bound}
\label{sec:tail-compression}
Keep the two disjoint bases of Section 16, possibly in different coordinates,
and write $p=15/7$, $s=5/3$. For either $H=Ph$ or $H=P\sqrt\pi$ let
$W_H=\pi/H^2$ and $T_H=M_{\sqrt{W_H}}P M_{\sqrt{W_H}}$.
$T_{Ph}$ is the exact symmetrized chain;
$T_{P\sqrt\pi}$ gives the unnormalized trial form. All norms below act on
counting probability space. The restriction to a set means killing outside it.

\paragraph{A failed candidate for a nonvanishing return.}
If $\|\pi\|_\infty\le B$, then for every syndrome set $E$,
\[
 \|1_E T_H1_E\|\le C_B\,\mu(E)^{2/5}.
\]
For any set of ambient words $F$, without requiring it to consist of
whole fibers, the corresponding bound is $C_B\mu(F)^{2/15}$.
Indeed Section 16 gives $\|H^{-1}\|_{10}\le C_B$ without matching,
hence $\|W_H\|_5\le C_B$. On the syndrome cube the square root of $P$
maps $L^2$ contractively into $L^5$ by the two basis steps, so the
Gram-product identity and H\"older bound the compression by
$\|1_E W_H\|_{5/3}$. Equivalently use (CR3) at its forward endpoint.
For an arbitrary ambient set the square root of ordinary half noise
maps $L^2$ into $L^3$, giving $\|1_F W_H\|_3$.
These give the stated exponents. The trial normalization costs at most $B^2$
as in Section 16. In particular the positive one-sided radial densities
of Section 17 have vanishing exact norm on \emph{every} sequence of
sets of vanishing uniform mass, although the trial inverse fifteenth
moment diverges. This does not give spectral recovery in the absence
of growing matching.

\paragraph{A stationary tail condition at the critical power.}
Suppose, for some fixed $a>0$ and $1/2<\gamma<1$, that
\[
 \sup_n \mathbb E_\mu \big[\pi^p
   \exp\{a(\log(e+\pi))^\gamma\}\big]<\infty .                 \tag{TB1}
\]
Both compressions tend uniformly to zero on syndrome sets with uniform
mass tending to zero, with no matching, rank bound, or permutation
invariance. Under the growing matching assumptions of (FF1) full
ordered spectral recovery follows at unrestricted rank.
The statement concerns arbitrary two-base quotients. No geometric
condition on their stationary level sets enters (TB1).

For a positive function $f$, let $f^*(u)$ be its decreasing quantile
under $\mu$, $0<u<1$, and put $L=\log(1/u)$. Section 16, using the
\href{https://www.cs.cmu.edu/~odonnell/papers/nicd.pdf}{binary reverse
hypercontractive inequality}, gives
$\log\|H^{-1}\|_{15-\delta}\le C_B/\delta$ for $0<\delta\le\delta_0$,
where $B$ bounds $\|\pi\|_p$. Markov's inequality with
$\delta=L^{-1/2}$ gives, for all sufficiently large $L$,
\[
 (H^{-1})^*(u)\le \exp\{L/15+C_B'\sqrt L\}.
\]
Write $y=\pi^*(u)$. Since $\mu(\pi\ge y)\ge u$, (TB1) gives
$p\log y+a(\log(e+y))^\gamma\le L+O(1)$. If $\log y\le L/(2p)$
there is already a linear saving. Otherwise the second term is at
least $a(L/(2p))^\gamma$. Thus for some $\eta>0$ and all large $L$,
$y\le\exp(L/p-\eta L^\gamma)$, uniformly across the sequence.

For a syndrome set $E$ of mass $m$, layer cake gives
\[
  \int_E \pi^s H^{-2s}\,d\mu
  \le \int_0^m (\pi^*(u))^s ((H^{-1})^*(u))^{2s}\,du .
\]
Indeed the mass of an intersection of three upper level sets is at most
the smallest of their masses. Using $s/p+2s/15=1$, the last integral
for all small $m$ is bounded by
\[
 \int_{\log(1/m)}^\infty
     \exp\{-s\eta L^\gamma+2sC_B'\sqrt L\}\,dL .
\]
It tends to zero since $\gamma>1/2$. Consequently $\|1_E W_H\|_s$
vanishes uniformly. The compression estimate above applies, including
on the ambient space for a set made of whole fibers by the positive
leading-vector reduction of (CR3). The exceptional multiplier set in
(FF1) has vanishing uniform mass by the strictly positive moment of
$W_{Ph}$ (alternatively by its uniform $L^s$ bound just proved).
The spectral assertion follows from (FF1).

\paragraph{Distinction from higher stationary powers.}
For $m$ syndrome bits with two identical bases, let $u=2^{-m}$ and
give one atom height $A_m=\exp(\log(1/u)/p-D(\log(1/u))^\gamma)$.
Take $\pi=(1+A_m 1_{\{0\}})/(1+u A_m)$.
It is normalized and strictly positive. For fixed $a,\gamma$ choose
$D>a/p^{1+\gamma}$. The contribution of the atom in (TB1) then
tends to zero, while $uA_m^{p+\epsilon}\to\infty$ for every
$\epsilon>0$. This distinguishes the stationary moment assumptions.
The example does not assert matching or establish sharpness of the
logarithmic correction.
